% =====================================================================
\section{Limitations}\label{sec:limitations}
% =====================================================================

This section enumerates the limitations of the analysis transparently.
The section is reorganised relative to the Mode-B-era version to put
the linearisation issue in its rightful place at the head, since it
materially weakened the original headline claim.

\subsection{Method-level limitations}

\begin{enumerate}
\item \textbf{Linearised analysis was misleading.} The earlier
  project iterations (the linearised stress test, the cross-dataset
  validation, and the cross-channel validation) used a Mode-B
  linearised response, where every angular observable is approximated
  as linear
  in $\Delta C_{9}$ via a central-difference slope at $\delta=\pm 0.5$
  around $C_{9}^{\mathrm{SM}}$. At the fitted shifts
  $|\Delta C_{9}|\sim 1$ this is wrong: per-bin linearisation
  residuals reach $4\sigma$ on LHCb 2025, and the headline
  $\Delta\mathrm{AIC}$ shifts from $-1.67$ (linearised) to $+1.09$
  (non-linear refit), reversing the kernel-vs-constant comparison.
  This paper now uses non-linear \texttt{flavio.np\_prediction}
  evaluation throughout for the headline numbers; Mode-B values are
  retained only as a methodology diagnostic. Future analyses on this
  anomaly should not use a linearised Taylor expansion.

\item \textbf{Variant choice is geometry-confirmed but data-discovered.}
  The unweighted Laplacian was originally selected because it gave
  the best LHCb-data $\chi^{2}$ among the three admissible variants
  (\S\ref{sec:derivation}, Table~\ref{tab:variant_selection}).
  We subsequently verified that the same variant wins on a pure
  geometry criterion (correlation with the Layer-1 continuum kernel,
  no LHCb data), and we now defend the AIC counting $k=1$ for VFD on
  this basis. The variant choice therefore does not consume an
  effective fit parameter, but the historical contingency that the
  data was looked at first is acknowledged.

\item \textbf{Single Wilson coefficient.} We compare against
  $\mathrm{FREE\_C9}$ and $\mathrm{FREE\_C9 + FREE\_C10}$ but do not
  test against the wider Wilson-coefficient space ($C_{9}'$, $C_{10}'$,
  scalar/tensor operators, lepton-flavour-non-universal shifts).
  The Mode-B stress test (\S\ref{sec:stress}) ruled out $C_{10}$ at
  the relevant level on the joint LHCb 2025 fit; the broader
  Wilson-coefficient comparison is deferred.

\item \textbf{Source ansatz.} The equatorial inner-product shell
  (30 vertices) is chosen as the source $f$ in
  Eq.~\eqref{eq:psi_discrete}. This is a natural choice for the
  2I-equivariant graph but is not derived from a first-principles
  closure operator on the $B\to K^{*}$ kinematic substrate; alternative
  sources have not been comprehensively scanned.

\item \textbf{Theoretical derivation incomplete.} The kernel is built
  from the standard graph Laplacian + $\varphi^{-2}$ regulariser plus
  the equatorial-shell source ansatz. The proper Hodge edge Laplacian
  $L_{\mathrm{edge}} = \delta_{2}\delta_{2}^{\mathsf T} +
  \delta_{1}^{\mathsf T}\delta_{1}$ with 2I-equivariant boundary maps
  from the 1200 triangular faces, and the Lyapunov-flow selection
  rule of \citep{SmartAdaptiveClosureTransport2026}, are deferred.
\end{enumerate}

\subsection{Data and statistical limitations}

\begin{enumerate}\setcounter{enumi}{5}
\item \textbf{Bootstrap covariance.} The bin bootstrap uses diagonal
  $\sigma_{\mathrm{stat}}^{2}+\sigma_{\mathrm{syst}}^{2}$ rather than
  the full LHCb covariance, because resampling with replacement
  produces duplicated rows and a singular covariance submatrix. This
  is the standard non-parametric bootstrap but it under-counts
  inter-bin correlation information that the all-data fit exploits.
  The bootstrap is therefore reported as a sign-stability test, not
  a confidence interval on the central-fit $A$.

\item \textbf{Form-factor MC scope.} The 20-sample BSZ form-factor
  Monte Carlo varies only the joint $B\to K^{*}$ form-factor
  multivariate constraint. CKM elements, quark masses, and decay
  constants are held at their flavio centrals. The 90\,\% CI
  $\pm 0.21$ on $A$ is a lower bound on the FF-induced uncertainty.
  The binomial 95\,\% CI on the ``$0/20$ negative'' result spans
  $[0, 14\,\%]$ — the form-factor MC is too small to give a tight
  constraint on the lower tail of the sign distribution.

\item \textbf{$B_{s}\to\phi$ covariance.} HEPData publishes no
  inter-observable correlation matrix for the LHCb 2015
  $B_{s}\to\phi$ angular fit. We use diagonal stat$\oplus$syst, which
  may overestimate the per-observable independence and therefore
  admit a larger best-fit amplitude than a covariance-aware fit
  would.

\item \textbf{High-$q^{2}$ flavio QCDF.} flavio's QCDF interpolation
  issues a ``do not trust above $6\,\mathrm{GeV}^{2}$'' warning. We
  use the values anyway. For bins above $6\,\mathrm{GeV}^{2}$ the SM
  prediction has theoretical uncertainty larger than the LHCb stat
  + syst errors; this is not currently propagated as a separate
  covariance source.

\item \textbf{Bound pinning.} Two of the cross-channel
  $B_{s}\!\to\!\phi$ region-split fits (central and high $q^{2}$)
  pinned at the optimiser bounds because slopes drop to
  $O(10^{-4})$ there. These are reported but flagged in the run log;
  only the unpinned low-$q^{2}$ region split is taken as quantitative.

\item \textbf{$P_{4}'$ convention.} CMS publishes $P_{4}'$ values
  outside the $[-1, +1]$ flavio-convention range. The full
  $P_{4}'$-included fit produces an inflated $\chi^{2}$. The
  CMS-no-$P_{4}'$ fit gives a clean result; it is the one used in
  the headline. The $P_{4}'$-included row is reported as a
  diagnostic of the convention mismatch.

\item \textbf{Single-observable basis on $B_{s}\to\phi$.} The 2015
  LHCb publication chose the $S$-basis, not the amplified $P$-basis.
  Future $B_{s}\to\phi$ analyses with $P$-basis output would tighten
  the cross-channel amplitude constraint substantially.
\end{enumerate}

\subsection{What is not in scope}

\begin{enumerate}\setcounter{enumi}{12}
\item \textbf{Bayesian posterior.} All intervals are frequentist
  ($\Delta\chi^{2}=1$ Gaussian approximation or non-parametric
  bootstrap). A proper Bayesian treatment marginalising over form
  factors and charm-loop nuisance parameters is not in scope.

\item \textbf{$B^{+}\to K^{+}\mu\mu$.} HEPData submissions for the
  relevant LHCb (1209.4284, 1403.8045, 1804.07167) and CMS PRD 2018
  papers are not publicly available. This decay channel is
  uncovered.

\item \textbf{Branching-fraction observables.} flavio's
  \texttt{<dBR/dq2>} returns $0$ in the version used; rate observables
  are uncovered.

\item \textbf{Belle / Belle II.} Belle's full $\Upsilon(4S)$ angular
  analysis lacks a public correlation matrix on HEPData; Belle II had
  not released a comparable angular dataset at the time of writing.
\end{enumerate}

\subsection{The honest verdict}

The result is consistency-of-direction, not statistical preference.
The geometry-derived kernel describes the angular anomaly with one
parameter per dataset, sign-uniformly, across two collaborations,
two isospin partners, and two decay channels. Under non-linear
evaluation it does \emph{not} beat a constant Wilson-coefficient shift
on AIC; the stacked Akaike weight is $\approx 1.9{:}1$ in favour of
$\mathrm{FREE\_C9}$. The substantive surviving claim is that
\emph{a single geometry-derived $q^{2}$-shape with no fitted shape
parameters and one fitted amplitude per dataset is sign-consistent
with the $b\to s\mu\mu$ angular anomaly across five public fits}, with
the stronger ``it beats the constant shift on AIC'' claim of the
earlier project iterations being a linearisation artifact.
